% year: תשפ"ו
% semester: א
% moed: א
% number: 2א
% points: 12
% categories: func-analysis, func-limits
% official_solution: yes
% source: exams/מבחנים/תשפו/AnyScanner_02_01_2026(1) (2).pdf

\begin{question}
(12 נק') חשבו את כל האסימפטוטות של הפונקציה
\[ f(x) = \frac{x^2 - 2\ln(x)}{x - 1} \]
\end{question}

\begin{hint}
התחילו מתחום ההגדרה: הנקודות ה"חשודות" לאסימפטוטה אנכית הן קצוות התחום ונקודות שבהן המכנה מתאפס. אסימפטוטה משופעת $y = ax + b$ יש לבדוק רק בכיוון שבו התחום אינסופי.
\end{hint}

\begin{solution}
\textbf{תחום הגדרה:} $\ln x$ מוגדר עבור $x > 0$, והמכנה מתאפס ב-$x = 1$. לכן התחום הוא $(0,1)\cup(1,\infty)$, ו-$f$ רציפה בו.

\textbf{אסימפטוטות אנכיות.} הנקודות החשודות הן $x = 0$ (קצה התחום) ו-$x = 1$.
\begin{itemize}
  \item $x \to 0^+$: המונה $x^2 - 2\ln x \to 0 - (-\infty) = +\infty$ והמכנה $x - 1 \to -1$, ולכן
  \[ \lim_{x\to 0^+} f(x) = -\infty. \]
  לכן $x = 0$ אסימפטוטה אנכית (מימין).
  \item $x \to 1$: המונה $\to 1 - 2\ln 1 = 1 > 0$ והמכנה $\to 0$, כאשר $x - 1 < 0$ משמאל ו-$x - 1 > 0$ מימין. לכן
  \[ \lim_{x\to 1^-} f(x) = -\infty, \qquad \lim_{x\to 1^+} f(x) = +\infty, \]
  ו-$x = 1$ אסימפטוטה אנכית (משני הצדדים).
\end{itemize}

\textbf{אסימפטוטות אופקיות/משופעות.} התחום אינסופי רק בכיוון $+\infty$, לכן בודקים רק $x \to +\infty$. נחפש $y = ax + b$:
\[ a = \lim_{x\to\infty} \frac{f(x)}{x} = \lim_{x\to\infty} \frac{x^2 - 2\ln x}{x^2 - x}
     = \lim_{x\to\infty} \frac{x^2\left(1 - \frac{2\ln x}{x^2}\right)}{x^2\left(1 - \frac{1}{x}\right)} = \frac{1 - 0}{1 - 0} = 1, \]
כי $\frac{\ln x}{x^2} \to 0$ (למשל לפי לופיטל: $\lim \frac{1/x}{2x} = 0$).
\[ b = \lim_{x\to\infty} \bigl(f(x) - x\bigr) = \lim_{x\to\infty} \frac{x^2 - 2\ln x - x^2 + x}{x - 1}
     = \lim_{x\to\infty} \frac{x\left(1 - \frac{2\ln x}{x}\right)}{x\left(1 - \frac{1}{x}\right)} = 1, \]
כי $\frac{\ln x}{x} \to 0$. מכיוון ששני הגבולות סופיים, הישר $y = x + 1$ הוא אסימפטוטה משופעת ב-$+\infty$ (ואין אסימפטוטה אופקית).

\textbf{תשובה:} אסימפטוטות אנכיות $x = 0$ ו-$x = 1$; אסימפטוטה משופעת $y = x + 1$ כאשר $x \to +\infty$.
\end{solution}
